\chapter{\ifproject%
\ifenglish Project Structure and Methodology\else โครงสร้างและขั้นตอนการทำงาน\fi
\else%
\ifenglish Project Structure\else โครงสร้างของโครงงาน\fi
\fi
}

บทนี้นำเสนอรายละเอียดของโครงสร้างระบบและการออกแบบของเกมที่พัฒนาขึ้น โดยอธิบายถึงองค์ประกอบหลักของระบบ และความสัมพันธ์ระหว่างแต่ละส่วน ซึ่งถูกออกแบบโดยอ้างอิงจากหลักการและทฤษฎีที่กล่าวไว้ในบทก่อนหน้า

\section{Core Gameplay System}

Core Gameplay System เป็นโครงสร้างหลักของเกมที่กำหนดรูปแบบการเล่นและลำดับของกิจกรรมที่ผู้เล่นต้องทำ โดยในโครงงานนี้ได้ออกแบบระบบให้มีลักษณะเป็นลูป (loop-based system) ซึ่งเชื่อมโยงระบบต่าง ๆ เข้าด้วยกันอย่างเป็นระบบ

Core Gameplay Loop ของเกมสามารถสรุปได้ดังนี้:

ล่ามอนสเตอร์ $\rightarrow$ ได้วัตถุดิบ $\rightarrow$ ทำอาหาร $\rightarrow$ แข่งขัน $\rightarrow$ พัฒนาตัวละคร $\rightarrow$ วนซ้ำ

\vspace{0.5em}
\noindent การออกแบบลูปดังกล่าวทำให้ผู้เล่นมีเป้าหมายที่ชัดเจนในแต่ละช่วงของการเล่น และสามารถเรียนรู้ผ่านการเล่นซ้ำ (learning by repetition)

\subsection{เมือง (Hub System)}

เมืองทำหน้าที่เป็นศูนย์กลางของเกม (safe zone) ที่ผู้เล่นสามารถเตรียมตัวก่อนออกผจญภัย โดยผู้เล่น \par

\noindent สามารถ:

\begin{itemize}
\item พูดคุยกับตัวละคร (NPC)
\item รับโจทย์การแข่งขัน
\item จัดการวัตถุดิบและเมนูอาหาร
\item ทำอาหารเพื่อเตรียมบัฟ
\end{itemize}
ระบบนี้ช่วยให้ผู้เล่นสามารถวางแผนและเลือกแนวทางการเล่นได้


\subsection{การออกสำรวจและต่อสู้ (Exploration \& Combat)}

ผู้เล่นจะออกไปยังพื้นที่ต่าง ๆ เช่น ป่า หรือดันเจี้ยน เพื่อทำการต่อสู้กับมอนสเตอร์และเก็บวัตถุดิบ โดยมอนสเตอร์แต่ละประเภทมีพฤติกรรมและรูปแบบการโจมตีที่แตกต่างกัน
\par\vspace{1em}\noindent
การออกแบบส่วนนี้ช่วยสร้างความท้าทายและเพิ่มความหลากหลายของ gameplay


\subsection{ระบบการทำอาหาร (Cooking System)}

หลังจากได้วัตถุดิบ ผู้เล่นสามารถนำมาปรุงอาหาร โดยระบบการทำอาหารถูกออกแบบให้เข้าใจง่าย และมีทั้งรูปแบบอัตโนมัติและมินิเกม 
\par\vspace{1em}\noindent
อาหารที่ได้จะส่งผลต่อความสามารถของตัวละคร เช่น การฟื้นฟูพลัง หรือการปลดล็อกสกิล

\subsection{ระบบการแข่งขัน (Competition System)}

ผู้เล่นสามารถนำอาหารที่ปรุงไปแข่งขัน โดยระบบจะมีโจทย์กำหนด เช่น วัตถุดิบและรสชาติเป้าหมาย
\par\vspace{1em}\noindent
ผลการแข่งขันจะถูกคำนวณจากระบบคะแนน ซึ่งรวมทั้งความถูกต้องของรสชาติและทักษะของผู้เล่น

\subsection{ระบบความก้าวหน้า (Progression System)}

เมื่อผู้เล่นชนะการแข่งขัน จะสามารถปลดล็อกพื้นที่ใหม่ มอนสเตอร์ใหม่ และเนื้อเรื่องเพิ่มเติม
\par\vspace{1em}\noindent
ระบบนี้ช่วยสร้างแรงจูงใจให้ผู้เล่นเล่นซ้ำและพัฒนาทักษะของตนเอง
\par\vspace{1em}\noindent
โดยสรุป Core Gameplay System ของเกมนี้เป็นการเชื่อมโยงระบบต่าง ๆ เข้าด้วยกันอย่างเป็นลำดับ \par
\noindent ทำให้เกิดการเล่นที่ต่อเนื่อง มีเป้าหมายชัดเจน และสนับสนุนการเรียนรู้ผ่านประสบการณ์การเล่น

\section{Food System}

ระบบอาหาร (Food System) เป็นระบบหลักของเกมที่ทำหน้าที่เชื่อมโยงการเล่นในหลายมิติทั้งด้านการต่อสู้ การพัฒนาตัวละคร และการดำเนินเนื้อเรื่อง โดยอาหารในเกมไม่ได้มีบทบาทเพียงแค่เป็นไอเทมสำหรับฟื้นฟูพลัง แต่ยังเป็นกลไกสำคัญในการกำหนดรูปแบบการเล่นของผู้เล่น

\subsection{บทบาทของอาหารในเกม}
อาหารในเกมนี้มีบทบาทหลัก 3 ด้าน ได้แก่
\begin{itemize}
\item \textbf{ไอเทม (Item)} ใช้สำหรับฟื้นฟูพลังชีวิต หรือเพิ่มสถานะ (buff) ให้กับตัวละคร
\item \textbf{สกิล (Skill Unlock Mechanism)} อาหารบางประเภทสามารถปลดล็อกสกิลใหม่ซึ่งส่งผลต่อรูปแบบการต่อสู้ของผู้เล่น
\item \textbf{ระบบพัฒนา (Progression System)} การค้นพบและปรุงอาหารใหม่ช่วยให้ผู้เล่นสามารถพัฒนาตัวละคร และปลดล็อกเนื้อหาเพิ่มเติม
\end{itemize}

\subsection{ประเภทของอาหาร}
ระบบอาหารแบ่งออกเป็น 2 ประเภทหลัก ได้แก่
\begin{enumerate}
\item \textbf{อาหารทั่วไป (Consume Food)} ใช้สำหรับการบริโภคเพื่อเพิ่มพลังหรือสถานะ โดยมีขั้นตอนการทำที่เรียบง่าย และสามารถทำซ้ำได้
\item \textbf{อาหารแข่งขัน (Competition Food)} ใช้ในระบบการแข่งขัน โดยมีโจทย์กำหนด เช่น วัตถุดิบและรสชาติเป้าหมาย ซึ่งต้องอาศัยทั้งการวางแผนและทักษะของผู้เล่น
\end{enumerate}

\subsection{กระบวนการทำอาหาร (Cooking Process)}

กระบวนการทำอาหารประกอบด้วย 2 ขั้นตอนหลัก

\begin{enumerate}
\item \textbf{การเลือกเมนู}  
ผู้เล่นสามารถเลือกเมนูจากรายการที่ปลดล็อกแล้ว หากมีวัตถุดิบครบสามารถทำได้ทันที

\item \textbf{การปรุงอาหาร}  
มี 2 รูปแบบ ได้แก่
\begin{itemize}
\item แบบอัตโนมัติ (Auto) สำหรับเมนูที่ผู้เล่นมีความชำนาญ
\item แบบมินิเกม (Mini Interaction) เช่น การกดจังหวะ การควบคุมระดับไฟ หรือการจัดลำดับขั้นตอน
\end{itemize}
\end{enumerate}

คุณภาพของอาหารจะขึ้นอยู่กับความแม่นยำของผู้เล่นในขั้นตอนนี้

\subsection{Food-to-Skill Transformation}

อาหารบางประเภทสามารถเปลี่ยนเป็นสกิลที่ใช้ในการต่อสู้ได้ โดยเมื่อผู้เล่นใช้สกิลจากอาหาร จะเกิดการ
\par\vspace{0.5em}\noindent
เปลี่ยนแปลงรูปแบบการต่อสู้ (Combat Mode) เช่น

\begin{itemize}
\item การโจมตีระยะใกล้ (Melee)
\item การโจมตีระยะไกล (Ranged)
\item การโจมตีที่มีสถานะพิเศษ (Status Effect)
\end{itemize}

ระบบนี้ช่วยเพิ่มความหลากหลายของ gameplay และเปิดโอกาสให้ผู้เล่นสามารถเลือกแนวทางการเล่น (playstyle) ได้

\subsection{ระบบความชำนาญ (Food Mastery System)}

ผู้เล่นสามารถพัฒนาความชำนาญของเมนูอาหารได้จากการทำซ้ำ ซึ่งจะส่งผลให้:

\begin{itemize}
\item สามารถทำอาหารได้รวดเร็วขึ้น
\item เพิ่มคุณภาพของอาหาร
\item ปลดล็อกตัวเลือกการทำอาหารแบบอัตโนมัติ
\end{itemize}

ระบบนี้ช่วยกระตุ้นให้ผู้เล่นมี “เมนูประจำตัว” และสร้างเอกลักษณ์ในการเล่น

\subsection{การเชื่อมโยงกับระบบอื่น}

Food System มีบทบาทสำคัญในการเชื่อมโยงกับระบบอื่นในเกม ได้แก่

\begin{itemize}
\item เชื่อมกับ Combat System ผ่านการปลดล็อกสกิล
\item เชื่อมกับ Competition System ผ่านการใช้เมนูในการแข่งขัน
\item เชื่อมกับ Progression System ผ่านการปลดล็อกเนื้อหาใหม่
\end{itemize}
\par\vspace{0.5em}\noindent
โดยสรุป ระบบอาหารเป็นแกนหลักของเกมที่เชื่อมโยงระบบต่าง ๆ เข้าด้วยกัน และเป็นปัจจัยสำคัญที่กำหนดทั้งรูปแบบการเล่นและความก้าวหน้าของผู้เล่น

\section{Combat System}

ระบบการต่อสู้ (Combat System) เป็นองค์ประกอบสำคัญของเกมที่ทำหน้าที่สร้างความท้าทายและความหลากหลายของการเล่น โดยในโครงงานนี้ ระบบการต่อสู้ถูกออกแบบให้เชื่อมโยงโดยตรงกับระบบอาหาร 
\par\vspace{0.5em}\noindent
(Food System) เพื่อให้เกิดความสัมพันธ์ระหว่างการเตรียมตัวและการต่อสู้

\subsection{รูปแบบการต่อสู้}

การต่อสู้ในเกมเป็นแบบเรียลไทม์ (real-time combat) โดยผู้เล่นสามารถเคลื่อนที่ หลบหลีก และโจมตีได้อย่างอิสระ ซึ่งช่วยให้เกิดความลื่นไหลของการเล่น และต้องอาศัยทักษะในการควบคุม
\par\vspace{0.5em}\noindent
ผู้เล่นต้องวิเคราะห์พฤติกรรมของมอนสเตอร์และเลือกจังหวะในการโจมตีหรือหลบหลีกอย่างเหมาะสม

\subsection{Combat Mode System}

ระบบ Combat Mode เป็นกลไกสำคัญที่เชื่อมระหว่างอาหารและการต่อสู้ โดยเมื่อผู้เล่นใช้สกิลจากอาหาร รูปแบบการต่อสู้จะเปลี่ยนไปตามประเภทของอาหารนั้น

ตัวอย่าง Combat Mode ได้แก่:

\begin{itemize}
\item \textbf{Melee Mode}  
การโจมตีระยะใกล้ มีความเร็วสูงและความเสียหายสมดุล

\item \textbf{Ranged Mode}  
การโจมตีระยะไกล เช่น ยิงลูกศรหรือพลังงาน แต่มีความเร็วในการโจมตีต่ำกว่า

\item \textbf{Special Mode}  
การโจมตีที่มีผลพิเศษ เช่น ทำให้ศัตรูติดสถานะ (stun, slow)
\end{itemize}

ระบบนี้ช่วยให้ผู้เล่นสามารถปรับเปลี่ยนรูปแบบการเล่น (playstyle) ได้ตามสถานการณ์

\subsection{Elemental Weakness System}

มอนสเตอร์แต่ละประเภทมีจุดอ่อน (elemental weakness) ที่แตกต่างกัน เช่น

\begin{itemize}
\item มอนสเตอร์ธาตุไฟ แพ้การโจมตีธาตุน้ำ
\item มอนสเตอร์ประเภทพืช แพ้ธาตุไฟ
\end{itemize}

ผู้เล่นจึงต้องเลือกสกิลจากอาหารให้เหมาะสมกับศัตรูในแต่ละด่าน ซึ่งเพิ่มมิติของการวางแผน

\subsection{Enemy Behavior Interaction}

พฤติกรรมของมอนสเตอร์ถูกออกแบบให้มีความหลากหลาย เช่น

\begin{itemize}
\item การโจมตีระยะใกล้
\item การยิงโปรเจกไทล์
\item การสร้างพื้นที่อันตราย (hazard zone)
\end{itemize}

ผู้เล่นต้องเรียนรู้รูปแบบการเคลื่อนไหวและการโจมตีของมอนสเตอร์เพื่อปรับกลยุทธ์ในการต่อสู้

\subsection{Combat Strategy Loop}

ระบบการต่อสู้เชื่อมโยงกับระบบอาหารในรูปแบบลูป ดังนี้:

ล่ามอนสเตอร์ $\rightarrow$ ได้วัตถุดิบ $\rightarrow$ ทำอาหาร $\rightarrow$ ปลดล็อกสกิล $\rightarrow$ เลือกสกิล $\rightarrow$ ต่อสู้
\par\vspace{0.5em}\noindent
ลูปนี้ทำให้การต่อสู้ไม่ใช่เพียงการใช้ทักษะ แต่ยังเกี่ยวข้องกับการวางแผนล่วงหน้าและการเลือกใช้ทรัพยากรอย่างเหมาะสม
\par\vspace{0.5em}\noindent
โดยสรุป ระบบการต่อสู้ในเกมนี้ถูกออกแบบให้มีความยืดหยุ่นและเชื่อมโยงกับระบบอาหารอย่างใกล้ชิด ทำให้เกิดความหลากหลายของรูปแบบการเล่น และเพิ่มความลึกของ gameplay

\section{Competition System}

ระบบการแข่งขันทำอาหาร (Competition System) เป็นกลไกสำคัญที่ใช้ในการวัดความสามารถของผู้เล่น และเป็นเงื่อนไขในการปลดล็อกเนื้อหาใหม่ภายในเกม โดยระบบนี้ผสมผสานทั้งการวางแผนล่วงหน้า (strategic planning) และทักษะในการเล่น (skill-based gameplay)

\subsection{โครงสร้างของการแข่งขัน}

การแข่งขันแต่ละครั้งจะมี “โจทย์” (Judge Requirement) ที่กำหนดโดยระบบ ซึ่งประกอบด้วย:

\begin{itemize}
\item เมนูอาหารที่ต้องทำ
\item วัตถุดิบหลัก (Required Ingredients)
\item รสชาติเป้าหมาย (Target Flavor)
\end{itemize}

ผู้เล่นจะต้องเตรียมวัตถุดิบจากการล่ามอนสเตอร์ และนำมาปรุงอาหารให้ตรงกับเงื่อนไขที่กำหนด

\subsection{ระบบการให้คะแนน (Scoring System)}

คะแนนของการแข่งขันถูกคำนวณจากหลายปัจจัย ได้แก่

\begin{itemize}
\item \textbf{Flavor Score}  
คำนวณจากความใกล้เคียงของรสชาติ โดยใช้ Manhattan Distance

\item \textbf{Cooking Skill Score}  
ได้จากความแม่นยำในการทำมินิเกม เช่น การกดจังหวะ และลำดับขั้นตอน

\item \textbf{Ingredient Bonus}  
ได้จากการใช้วัตถุดิบที่มีความหายาก
\end{itemize}

คะแนนรวมสามารถเขียนได้เป็น:
\par
Final Score = Flavor Score + Cooking Skill Score + Ingredient Bonus

\subsection{มินิเกมทำอาหาร (Cooking Mini-game)}

ในระหว่างการปรุงอาหาร ผู้เล่นจะต้องทำมินิเกม เช่น
\begin{itemize}
\item กดปุ่มตามจังหวะ
\item ควบคุมระดับไฟ
\item จัดลำดับขั้นตอน
\end{itemize}
ความแม่นยำของผู้เล่นจะส่งผลต่อคุณภาพของอาหารและคะแนนที่ได้รับ

\subsection{ระบบตัดสินกรณีคะแนนเท่ากัน (Tie-Breaking)}

ในกรณีที่ผู้เล่นและคู่แข่งมีคะแนนรวมเท่ากัน ระบบจะใช้เกณฑ์เพิ่มเติมในการตัดสิน ได้แก่:

\begin{enumerate}
\item  Cooking Skill Score (ผู้ที่มีทักษะสูงกว่าจะชนะ)
\item ความแม่นยำของรสชาติ (Flavor Precision)
\item ความหายากของวัตถุดิบ (Ingredient Rarity)
\item เวลาในการทำอาหาร (Time Efficiency)
\end{enumerate}
หากยังไม่สามารถตัดสินได้ ระบบจะสุ่มผู้ชนะเพื่อป้องกันการเกิดผลลัพธ์ซ้ำ

\subsection{บทบาทของระบบการแข่งขันในเกม}

ระบบการแข่งขันทำหน้าที่เป็น ``Progression Gate'' ที่ผู้เล่นต้องผ่านเพื่อปลดล็อกพื้นที่ใหม่ มอนสเตอร์ใหม่ และเนื้อเรื่องเพิ่มเติม

\par\vspace{0.5em}\noindent
นอกจากนี้ ระบบยังช่วยเชื่อมโยงการเล่นในหลายมิติ ได้แก่

\begin{itemize}
\item การล่ามอนสเตอร์เพื่อเตรียมวัตถุดิบ
\item การวางแผนเลือกเมนู
\item การใช้ทักษะในการทำมินิเกม
\end{itemize}
\par\vspace{0.5em}\noindent
โดยสรุป ระบบการแข่งขันทำหน้าที่เป็นจุดทดสอบความสามารถของผู้เล่น และเป็นกลไกสำคัญในการผลักดันความก้าวหน้าในเกม

\section{System Integration}

ระบบต่าง ๆ ภายในเกมถูกออกแบบให้ทำงานร่วมกันอย่างเป็นระบบ (integrated system) โดยมีการเชื่อมโยงข้อมูลและการทำงานระหว่างแต่ละระบบ เพื่อสร้างประสบการณ์การเล่นที่ต่อเนื่องและมีความหมาย

\subsection{ภาพรวมการเชื่อมโยงระบบ}

ระบบหลักของเกมประกอบด้วย:
\begin{itemize}
    \item Combat System
    \item Food System
    \item Competition System
    \item Behavior System
    \item Progression System
\end{itemize}

ระบบเหล่านี้ไม่ได้ทำงานแยกกัน แต่ถูกเชื่อมโยงผ่าน Core Gameplay Loop ทำให้การกระทำในแต่ละระบบส่งผลต่อระบบอื่นโดยตรง

\subsection{การไหลของระบบ (System Flow)}

การทำงานของระบบสามารถอธิบายเป็นลำดับได้ดังนี้:
\par\vspace{0.5em}\noindent
ล่ามอนสเตอร์ $\rightarrow$ ได้วัตถุดิบ $\rightarrow$ ทำอาหาร $\rightarrow$ ปลดล็อกสกิล $\rightarrow$ เลือกสกิล $\rightarrow$ ต่อสู้ $\rightarrow$ แข่งขัน$\rightarrow$ ปลดล็อกเนื้อหาใหม่

ในลำดับดังกล่าว:
\begin{itemize}
    \item Combat System เป็นแหล่งของวัตถุดิบ
    \item Food System ใช้ในการแปลงวัตถุดิบเป็นความสามารถ
    \item Competition System ใช้ประเมินผลและควบคุม progression
\end{itemize}

\subsection{การเชื่อมโยงข้อมูลระหว่างระบบ}

ข้อมูลสำคัญที่ถูกส่งต่อระหว่างระบบ ได้แก่:
\begin{itemize}
    \item วัตถุดิบ (Ingredients) \\
    ได้จาก Combat และถูกใช้ใน Food System
    \item อาหาร (Food Data) \\
    ถูกใช้ทั้งใน Combat (สกิล) และ Competition (คะแนน)
    \item สกิล (Skill Data) \\
    ส่งผลต่อรูปแบบการต่อสู้ของผู้เล่น
    \item คะแนน (Score Data) \\
    ใช้ในการตัดสินผลการแข่งขันและ progression
\end{itemize}

\subsection{ความสัมพันธ์ระหว่างระบบ}

ระบบภายในเกมมีความสัมพันธ์แบบ dependency ซึ่งแต่ละระบบส่งผลต่อกัน เช่น:
\begin{itemize}
    \item Combat $\rightarrow$ ส่งผลต่อ Food (วัตถุดิบ)
    \item Food $\rightarrow$ ส่งผลต่อ Combat (สกิล)
    \item Competition $\rightarrow$ ส่งผลต่อ Progression
    \item Behavior System $\rightarrow$ ส่งผลต่อความยากและประสบการณ์การเล่น
\end{itemize}

ความสัมพันธ์นี้ทำให้ทุกระบบมีบทบาทสำคัญ และไม่มีระบบใดที่ทำงานอย่างโดดเดี่ยว

\subsection{ประโยชน์ของการออกแบบแบบ Integrated System}

การออกแบบระบบให้เชื่อมโยงกันช่วยให้:
\begin{itemize}
    \item เกมมีความลื่นไหล (flow) และต่อเนื่อง
    \item ผู้เล่นเข้าใจความสัมพันธ์ของระบบได้ง่าย
    \item เพิ่มความลึกของ gameplay
    \item สนับสนุนการเล่นซ้ำ
\end{itemize}


โดยสรุป การออกแบบ System Integration ในโครงงานนี้ช่วยให้ระบบต่าง ๆ ทำงานร่วมกันอย่างมีประสิทธิภาพ และสร้างประสบการณ์การเล่นที่เป็นหนึ่งเดียว